\subsection{Non-Convex Solvers}\label{app:non-convex-solvers}

\begin{table*}
	\centering
	\caption{Test accuracies for convex and non-convex formulations of the Gated ReLU training problem on a subset of 20 datasets selected from the UCI dataset repository
		Results are shown as \texttt{median (first-quartile/third-quartile)} for each method.}%
	\label{table:uci-gated-accuracies}
	\vspace{0.1in}
	\begin{small}
		\begin{tabular}{lccc} \toprule
			\textbf{Dataset}     & \textbf{C-GReLU} & \textbf{NC-GReLU (Adam)} & \textbf{NC-GReLU (SGD)} \\ \midrule
			magic                & 86.9 (86.8/87.0) & 82.9 (82.9/83.1)         & 82.1 (82.1/82.2)        \\
			statlog-heart        & 79.6 (79.6/79.6) & 85.2 (83.3/85.2)         & 83.3 (83.3/83.3)        \\
			mushroom             & 100.0 (100/100)  & 97.6 (97.6/97.9)         & 96.9 (96.9/96.9)        \\
			vertebral-column     & 87.1 (83.9/87.1) & 90.3 (90.3/91.9)         & 90.3 (90.3/90.3)        \\
			cardiotocography     & 90.1 (89.9/90.4) & 85.6 (85.6/85.9)         & 85.2 (85.2/85.4)        \\
			abalone              & 63.8 (63.7/64.1) & 58.7 (58.6/58.7)         & 58.1 (58.1/58.1)        \\
			annealing            & 90.6 (90.6/91.2) & 86.2 (86.2/86.8)         & 86.2 (85.5/86.2)        \\
			car                  & 89.9 (89.9/90.1) & 83.8 (83.8/84.1)         & 83.2 (82.9/83.2)        \\
			bank                 & 89.8 (89.7/89.9) & 89.9 (89.9/90.0)         & 89.8 (89.8/90.0)        \\
			breast-cancer        & 68.4 (68.4/68.4) & 68.4 (68.4/70.2)         & 70.2 (70.2/70.2)        \\
			page-blocks          & 96.8 (96.8/96.9) & 92.1 (92.0/92.1)         & 92.4 (92.3/92.4)        \\
			contrac              & 45.9 (45.6/46.3) & 53.1 (53.1/53.1)         & 53.4 (53.1/53.7)        \\
			congressional-voting & 63.2 (63.2/63.2) & 64.4 (64.4/64.4)         & 66.7 (66.7/66.7)        \\
			spambase             & 93.4 (93.2/93.4) & 91.6 (91.6/91.6)         & 91.2 (91.2/91.3)        \\
			synthetic-control    & 97.5 (97.5/97.5) & 98.3 (98.3/98.3)         & 97.5 (97.5/98.3)        \\
			musk-1               & 93.7 (91.6/93.7) & 93.7 (93.7/93.7)         & 94.7 (92.6/94.7)        \\
			ringnorm             & 69.8 (69.5/69.9) & 77.0 (77.0/77.0)         & 77.2 (77.1/77.2)        \\
			ecoli                & 82.1 (82.1/82.1) & 79.1 (79.1/80.6)         & 4.5 (3.0/43.3)          \\
			monks-2              & 69.7 (66.7/69.7) & 66.7 (66.7/66.7)         & 60.6 (57.6/63.6)        \\
			hill-valley          & 62.0 (59.5/66.1) & 57.0 (55.4/57.9)         & 58.7 (58.7/59.5)        \\ \bottomrule
		\end{tabular}
	\end{small}
\end{table*}


This section gives experimental details and additional results for experiments comparing the generalization performance of our convex reformulations to neural networks trained by optimizing the non-convex loss with stochastic gradient methods.

\textbf{Experimental Details}:
We selected 20 datasets randomly from our filtered collection of 73 datasets; see Appendix~\ref{app:uci-datasets} for details how the 73 datasets were obtained.

We used the default parameters for each of the convex solvers as described in Appendix~\ref{app:default-parameters}.
R-FISTA was limited to \( 2000 \) iterations, while SGD and Adam were limited to \( 2000 \) epochs.
For SGD and Adam, considered step-sizes from the following grid: \( \cbr{10, 5, 1, 0.5, 0.1, 0.01, 0.001} \).
We used a ``step'' decrease schedule for the step-sizes, dividing them by \( 2 \) every 100 epochs, which we found to work much better than the classical Robbins-Monro schedule~\citep{robbins1951sgd}.
We considered the following grid of ten regularization parameters:
\( \{
1 \times 10^{-6},
2.78 \times 10^{-6},
7.74 \times 10^{-6},
2.15 \times 10^{-5},
5.99 \times 10^{-5},
1.67 \times 10^{-4},
4.64 \times 10^{-4},
1.29 \times 10^{-3}
3.59 \times 10^{-2},
1.0 \times 10^{-2} \}
\).
For each method-dataset pair, we performed five-fold cross validation on the training set and selected the best step-size and regularization parameter according to the cross-validated test accuracy.

For the gated ReLU problems (both C-GReLU and NC-GReLU) we sampled the same set of \( 5000 \) activation patterns for both the convex reformulation and the original non-convex model.
We used \( 2500 \) activation patterns for the C-ReLU problem.
To ensure a similar model space, we computed at the number of active neurons (e.g. \( v_i \neq 0 \) or \( w_i \neq 0 \)) at convergence for C-ReLU and then used this as the number of hidden units for \( NC-ReLU \) problems.
Note that we extend our convex reformulations to multi-class problems using the results in Appendix~\ref{app:multi-class-extension}.
Similarly, we use the vector-output variant of the NC-ReLU problem (Eq.~\ref{convex-forms:eq:vector-non-convex-relu-mlp}) for multi-class problems.

After selecting hyper-parameters, we obtain the final test accuracies by re-training on the full training set and testing on a held-out test set.
To control for noise in the sampling of gate vectors in the Gated ReLU problems and C-ReLU, we repeat this final testing procedure five times with different random seeds.

\begin{table*}
	\centering
	\caption{Test accuracies for convex and non-convex formulations of the ReLU training problem on a subset of 20 datasets selected from the UCI dataset repository
		Results are shown as \texttt{median (first-quartile/third-quartile)} for each method.}%
	\label{table:uci-relu-accuracies}
	\vspace{0.1in}
	\begin{small}
		\begin{tabular}{lccc} \toprule
			\textbf{Dataset}     & \textbf{C-ReLU}  & \textbf{NC-ReLU (Adam)} & \textbf{NC-ReLU (SGD)} \\ \midrule
			magic                & 85.9 (85.8/85.9) & 86.9 (86.9/86.9)        & 86.4 (86.3/86.4)       \\
			statlog-heart        & 83.3 (81.5/83.3) & 83.3 (83.3/83.3)        & 79.6 (79.6/79.6)       \\
			mushroom             & 100.0 (100/100)  & 100.0 (100/100)         & 99.9 (99.9/99.9)       \\
			vertebral-column     & 90.3 (88.7/90.3) & 90.3 (90.3/90.3)        & 88.7 (88.7/88.7)       \\
			cardiotocography     & 89.9 (89.9/89.9) & 36.5 (22.8/36.5)        & 88.9 (88.9/88.9)       \\
			abalone              & 66.2 (66.1/66.3) & 65.3 (64.9/65.4)        & 66.1 (66.1/66.1)       \\
			annealing            & 90.6 (89.9/90.6) & 93.7 (93.7/93.7)        & 88.7 (88.1/88.7)       \\
			car                  & 87.8 (87.8/87.8) & 94.8 (94.8/94.8)        & 90.1 (90.1/90.1)       \\
			bank                 & 89.8 (89.7/89.9) & 90.8 (90.8/90.9)        & 90.5 (90.5/90.5)       \\
			breast-cancer        & 68.4 (66.7/68.4) & 64.9 (64.9/64.9)        & 68.4 (68.4/68.4)       \\
			page-blocks          & 94.0 (94.0/94.0) & 97.1 (97.1/97.1)        & 96.9 (96.9/96.9)       \\
			contrac              & 55.1 (54.1/55.4) & 54.4 (54.1/54.4)        & 53.7 (53.7/53.7)       \\
			congressional-voting & 63.2 (63.2/65.5) & 62.1 (62.1/62.1)        & 67.8 (67.8/67.8)       \\
			spambase             & 93.3 (93.2/93.4) & 93.5 (93.5/93.5)        & 93.2 (93.2/93.2)       \\
			synthetic-control    & 98.3 (97.5/98.3) & 96.7 (96.7/96.7)        & 96.7 (96.7/96.7)       \\
			musk-1               & 93.7 (93.7/93.7) & 96.8 (96.8/96.8)        & 95.8 (95.8/95.8)       \\
			ringnorm             & 77.0 (76.8/77.0) & 77.3 (77.3/77.4)        & 77.4 (77.3/77.5)       \\
			ecoli                & 80.6 (80.6/80.6) & 82.1 (82.1/82.1)        & 80.6 (80.6/80.6)       \\
			monks-2              & 69.7 (69.7/72.7) & 69.7 (66.7/69.7)        & 72.7 (72.7/75.8)       \\
			hill-valley          & 65.3 (64.5/65.3) & 62.8 (62.8/62.8)        & 55.4 (55.4/55.4)       \\ \bottomrule
		\end{tabular}
	\end{small}
\end{table*}


\textbf{Additional Results}:
Tables~\ref{table:uci-gated-accuracies} and~\ref{table:uci-relu-accuracies} report median test accuracies as well as first and third quartiles for the convex and non-convex formulations with gated ReLU and ReLU activations, respectively.
Note that these results are identical to those the provided in the main paper (\cref{table:non-convex-solvers}) but for the inclusion of variance/distribution information in the form of quartiles.

